\section{从相互作用到引力场}
\subsection{研究对象、距离与质量}

两个可视为质点的物体相距 $r$，引力大小为 $Gm_1m_2/r^2$，方向沿两质点连线。这里的 $r$ 是中心间距离；只有物体尺寸远小于距离，或者已经利用球对称证明外场等效时，才可直接代入。对任意延展物体，应先把物体分成质量微元，按矢量叠加求场，不能把外形半径当成万能的距离修正。

把很小的试验质量 $m$ 放在给定位置，定义引力场强 $\bm g=\bm F/m$。在牛顿力学中采用惯性质量与引力质量相等的实验事实，同一位置的自由落体加速度便不依赖试验物体的质量。场强是空间位置的函数，力还要乘受力物体的质量；不能把两者的量纲混用。

若以场源球心为原点，径向单位矢量向外，外场为 $\bm g=-GM\bm r/r^3$。负号说明吸引，$GM/r^2$ 只给大小。多天体叠加时，先画出各个指向场源的箭头，再选坐标分解；标量大小直接相加只适用于方向相同的特殊位置。

\typeout{ASTROFIG:fig:at1:field}\fig{fig:at1:field}{场点与场源之间的距离决定引力方向和大小}{\begin{tikzpicture}[scale=.85]
\draw[world](0,0)circle(1.3);\fill(0,0)circle(1.5pt)node[below]{$O$};\draw[dim](0,0)--(35:1.3)node[midway,above]{$R$};\node at(-.65,.55){$M$};\fill(3.8,0)circle(2pt)node[above]{$m$};\draw[force](3.8,0)--(2.5,0)node[midway,above]{$\bm F$};\draw[dim](0,-.8)--(3.8,-.8)node[midway,below]{$r$};
\end{tikzpicture}}
图\ref{fig:at1:field}中，场点在球外，公式中的距离是从 $O$ 到 $m$ 的整段 $r$，并非表面到场点的高度。若高度记为 $h$，应先写 $r=R+h$。\subsection{引力势、势能和功的符号}

取无穷远势能为零，球外引力势为 $\Phi=-GM/r$，质量 $m$ 的势能为 $U=m\Phi$。物体由 $r_1$ 移到 $r_2$，引力做功
\[
 W=\int_{r_1}^{r_2}-\frac{GMm}{r^2}\dd r
 =GMm\left(\frac1{r_2}-\frac1{r_1}\right)=-\Delta U.
\]
向内运动时势能更负，引力做正功；向外运动时正好相反。若没有推力、阻力和碰撞，机械能守恒。不同参考系下的动能不同，使用能量式前必须说明场源是否固定、两天体是否都在运动。


\section{球对称分布的场与内部结构}
\subsection{球壳定理怎样减少积分}

在均匀薄球壳内任取一点 $P$，从 $P$ 作很窄的双锥。两端截出的面积随距离平方增加，因此两块质量之比等于距离平方之比；它们的引力大小相等、方向相反。将所有双锥相加，球壳内部合场为零。这个论证比较的是同一微小立体角截出的两块面元，不是把球壳简单分成质量相等的两半。

对球外点，球壳积分给出势 $-GM/r$，因此其外场与质量集中在球心相同。将一般球对称物体分成同心薄壳，半径 $r$ 以内的质量产生 $GM(r)/r^2$，外部薄壳在该处不产生场。于是
\[
 \bm g(r)=-\frac{GM(r)}{r^3}\bm r,\qquad
 M(r)=4\pi\int_0^r\rho(s)s^2\dd s.
\]
均匀实心球有 $M(r)=4\pi\rho r^3/3$，所以球内场与位移成正比，$\bm g=-4\pi G\rho\bm r/3$。球心处场为零，表面处与外场连续；过了表面以后，场强才按距离平方反比减小。

\typeout{ASTROFIG:fig:at1:shell}\fig{fig:at1:shell}{内球质量产生场，外侧球壳只改变势的常数}{\begin{tikzpicture}[scale=.85]
\draw[world](0,0)circle(1.3);\fill(0,0)circle(1.5pt)node[below]{$O$};\draw[dim](0,0)--(35:1.3)node[midway,above]{$R$};\node at(-.65,.55){$M(r)$};\draw[guide](0,0)circle(.7);\fill(.7,0)circle(1.8pt)node[above]{$P$};\node at(3,.7){$M(r)$：半径 $r$ 内};\node at(3,-.2){外壳的合场为零};
\end{tikzpicture}}
图\ref{fig:at1:shell}区分了“外壳不产生合场”和“外壳不产生电势式的标量贡献”：引力势仍有外壳贡献。因此球内的势不能直接写成 $-GM(r)/r$。\subsection{球壳外部等效的势积分}

  球外等效点质量也可直接证明。取半径 $a$、质量 $M_s$ 的均匀薄球壳，场点距球心 $r$。将壳面按极角 $\theta$ 切成细环，环到场点的距离均为 $s=\sqrt{r^2+a^2-2ar\cos\theta}$，质量 $\dd M=(M_s/2)\sin\theta\dd\theta$。同一细环的势可先合并，所以
  \[\Phi(r)=-\frac{GM_s}{2}\int_0^\pi\frac{\sin\theta\dd\theta}{\sqrt{r^2+a^2-2ar\cos\theta}}.
  \]
  由 $2s\dd s=2ar\sin\theta\dd\theta$，积分化为
  \[\Phi(r)=-\frac{GM_s}{2ar}\int_{|r-a|}^{r+a}\dd s
  =-\frac{GM_s}{2ar}\big[r+a-|r-a|\big].\]
  球外 $r>a$，括号为 $2a$，得到 $-GM_s/r$；球内 $r<a$，括号为 $2r$，得到常数 $-GM_s/a$。分别对 $r$ 求负梯度，就同时给出球外平方反比场和球内零场。球心处可直接用等距性求势，或者取上式的连续极限，不把中间式的 $1/r$ 当作真实发散。
  
\typeout{ASTROFIG:fig:at1:ringintegral}\fig{fig:at1:ringintegral}{同一细环上到轴线场点的距离相同，适合先累加势}{\begin{tikzpicture}[scale=.85]
\draw[world](0,0)circle(1.7);\draw[axis](-2,0)--(4,0);\fill(3.5,0)circle(2pt)node[above]{$P$};\draw[blue,very thick](.85,-1.472)--(.85,1.472);\draw[guide](0,0)--(.85,1.472)node[midway,left]{$a$};\draw[guide](3.5,0)--(.85,1.472)node[midway,above]{$s$};\node[below]at(1.7,0){$r$};
\end{tikzpicture}}

  图\ref{fig:at1:ringintegral}中蓝线表示细环在轴截面内的截迹，不是一根有质量的直杆。利用圆环轴对称性合并标量势，比逐个分解三维引力分量更简洁；但要保留两端距离 $|r-a|,r+a$，绝对值正是内外分段不同的原因。若忽略绝对值，球内也会错误得到平方反比场。

  这种“先叠加势、再求导”的方法适用于质量分布固定、积分和求导可以交换的区域。在有面密度的理想薄壳表面，势连续而径向导数可以跳变；不能从势连续推出场也连续。对实心连续密度球，表面不存在额外面质量，内外场才连续连接。两种边界的物理区别应在画图时注明。

  对挖空分布，用负密度叠加是计算工具，真实物质并没有负质量。必须先保证负球完全位于对应正球中，并在各自适用的内外场区域使用公式。偏心球腔内任一点同时位于完整球与被挖去小球内部，所以两个线性场恰好相消掉场点变量；若把场点移到小球外面，就不能继续宣称场仍均匀。
  \subsection{球内势与流体压强}

用表面势连续连接球内场，均匀球内部有
\[
 \Phi(r)=-\frac{GM}{R}-\int_R^r g_s(s)\dd s
 =-\frac{GM}{2R}\left(3-\frac{r^2}{R^2}\right),
\]
其中有向径向分量 $g_s=-GMs/R^3$。球心势最小而场为零，两者并不矛盾：场取决于势的空间导数。

静止、球对称且不自转的液态天体，薄层受到的压强差平衡向内重力，$\dd p/\dd r=-\rho(r)GM(r)/r^2$。必须同时计算内部质量 $M(r)$ 和压强；把表面 $g$ 当作整个星球内部的常数会高估中心压强。密度是否近似不变也是额外物性假设，并非流体静力学方程自动给出的结论。


\section{圆轨道与共同质心}
\subsection{向心方程的适用条件}

小质量卫星绕固定球形天体作半径 $r$ 的圆周运动，由径向方程
\[
 \frac{GMm}{r^2}=m\frac{v^2}{r}=m\omega^2r
\]
依次得到 $v=\sqrt{GM/r}$、$\omega=\sqrt{GM/r^3}$、$T=2\pi\sqrt{r^3/(GM)}$。它们只适用于圆轨道。对于椭圆上的一个位置，径向加速度还含径向速度的变化，不能再写成 $v^2/r$；曲率半径也通常不等于到焦点的距离。

已知 $r,T$ 可以反求 $GM$。若只测得周期而不知道轨道尺度，不能唯一确定中心质量。地面附近测得的有效重力还可能包括自转影响，只有忽略自转或经过修正后，才能用 $gR^2$ 代替 $GM$。

\subsection{两颗天体都运动时}

两天体间距离为 $r$，共同质心将连线分成 $r_1=m_2r/(m_1+m_2)$、$r_2=m_1r/(m_1+m_2)$。若两者作圆周运动，角速度相同，分别列向心方程便得 $\omega^2=G(m_1+m_2)/r^3$。

一般相对运动用 $\bm r=\bm r_1-\bm r_2$ 描述，满足 $\ddot{\bm r}=-G(m_1+m_2)\bm r/r^3$。两体总动能可以分成质心平动和相对运动：$K=M_{\Sigma}V_C^2/2+m_*\dot{\bm r}^{,2}/2$，其中 $m_*=m_1m_2/M_{\Sigma}$。相对运动的引力参数是 $GM_{\Sigma}$，相对动能的质量是约化质量 $m_*$；这两个“质量位置”不同，不能相互替换。


\section{自转、潮汐与微小差别}
\subsection{支持力和引力不总相等}

赤道上与星球一起转动的物体，向内引力减去向外支持力，提供向心力：$mg_0-N=m\Omega^2R$。秤示的有效重力为 $N/m=g_0-\Omega^2R$。临界支持力为零，只表示物体无需接触力也能随表面转动；不会使它立即获得逃逸速度。若表面物体恰达到圆轨道速度，它仍被引力束缚。

在其他纬度，自转向心方向指向转轴，不严格沿球心方向，支持力也可能有切向分量；本讲需要精确关系时会明确坐标和近似。由临界转速估计密度时，还假设天体近似球形、表面依靠引力束缚而无显著材料强度。

\subsection{潮汐应比较同一参考系中的加速度}

取随中心天体平移、坐标轴不转动的自由落体参考系。远处质量 $M'$ 距原点 $D$，沿 $+x$ 方向。先减去它对原点的共同加速度，再对小位移 $\bm r=(x,y,z)$ 展开，得
\[
 \Delta\bm g=\frac{GM'}{D^3}(2x,-y,-z)+O(r^2/D^4).
\]
沿连线方向伸长，垂直方向压缩。近侧和远侧都相对中心向外偏移，原因是场的差别，不是远侧被另一个引力“向外吸引”。若选择转动坐标系，还必须加入该坐标系实际角速度产生的离心项和科里奥利项；不能把不同参考系的系数拼进同一个潮汐式。

对 $h/R$、$x/D$ 这类小参数，先写出未展开的式子，再保留所需阶数。结果应同时检查量纲和极限：远处天体距离趋于无穷时潮汐消失，极小高度差应给线性的测量差别。


\section{分布、测量与引力能的进一步联系}
\subsection{密度变化怎样进入引力极值}

  对球对称但不均匀的天体，$g(r)=GM(r)/r^2$。求导时包围质量也随半径变，
  \[g'(r)=4\pi G\rho(r)-\frac{2GM(r)}{r^3}.\]
  因而内部场强的极值条件是当地密度 $\rho(r)$ 等于包围球平均密度的 $2/3$。这只是候选条件，极大还是极小仍需检查两侧符号。均匀球的当地密度始终等于平均密度，故内部场一直增大；密度向外下降足够快时，最大场可以在表面以内。

  反过来，若观测给定球内的 $g(r)$，由 $M(r)=r^2g(r)/G$ 可恢复 $\rho(r)=[2g/r+g']/(4\pi G)$。这是用场反推分布，求导会放大测量噪声，因此数学可恢复不等于实验上可以不经平滑就直接逐点求导。
  
\subsection{自引力能的装配积分}

  把球形物体从内向外装配。已有内部质量 $M(r)$，从无穷远搬入薄壳 $\dd M$ 到半径 $r$，缓慢外力功为 $-GM(r)\dd M/r$。每一对质量只在后来加入那一份时计算一次，因此
  \[U=-\int_0^R\frac{GM(r)}r\dd M(r).\]
  均匀球代入 $M(r)=Mr^3/R^3$、$\dd M=3Mr^2\dd r/R^3$，给 $U=-3GM^2/(5R)$。若已经逐层装配，就不额外乘 $1/2$；若改用 $\tfrac12\int\Phi\dd M$，势中含全部质量贡献，才需要一半避免重复计数。

  把天体完全分散到无穷远、且末态动能趋零，所需最少外力功为 $-U$。这个总体束缚能不同于把单个小试验物从表面搬到无穷远的 $GMm/R$。前者改变整个场源分布，后者假设场源基本不受影响。
  
\subsection{观测量与未知量的对应}

  周期与轨道尺度通常首先确定引力参数 $GM$。圆轨道双星的两径向速度振幅还与轨道倾角有关；没有倾角时只能确定包含 $\sin^3 i$ 的质量组合，不能把沿视线速度直接当成真实轨道速度。相同的代数式若忽略投影，会产生系统性的质量低估。

  近似的条件应与测量量相配。用地面落下时间求 $g$，要求局部位移小；由星表转速约束密度，要求天体近球形、表面主要靠引力束缚；从隧道周期反推平均密度，还要明确内部是否均匀。题目中的“忽略”是建立简化模型的条件，不能在求得结果后再忘掉这些限制。
  

